\clearpage
\SourceLeaf{247}%
\Needspace{7\baselineskip}
\begin{center}
{\Large[\textsc{Congruenz und Symmetrie.}]}
\end{center}

\par\bigskip

\Needspace{7\baselineskip}
\begin{center}
[\textsc{Gerling} \emph{an} \textsc{Gauss}.\quad \emph{Cassel, 25. März 1813.}]
\end{center}

\{\mbox{.\ .\ .\ .\ .\ .} Im Decemberheft der monatlichen Correspondenz vom vorigen Jahr
erklärt Professor \textsc{Mollweide} den gewöhnlichen Beweis für die
Oberfläche der sphärischen Dreiecke für unstatthaft, indem dabei zwei
Dreiecke als gleich gebraucht werden, die 3 gleiche Seiten haben, aber wegen
verschiedener Lage derselben nicht congruent sein können, dieser Satz aber
in den sphärischen Trigonometrien nicht vorkommt. Ich habe von diesem Satz
einen Beweis gefunden, der mir sehr einfach zu sein scheint. Es ist
folgender: Zwei solche Dreiecke entstehen immer, wenn man die das eine
bildenden Kugelradien, jenseits des Mittelpunkts, verlängert. Legt man nun
durch die Winkelpunkte beider Dreiecke Ebenen, so erhellt sehr leicht, dass
sie parallel sind, und daraus folgt dann gleichfalls sehr leicht, dass sie
gleich weit vom Mittelpunkte abstehen. --- Sie beschreiben also auf der
Kugel congruente kleine Kreise, und congruente Oberflächensegmente, und die
Dreiecke entstehen, wenn man von diesen congruenten Oberflächensegmenten die
3 Zweiecke abzieht, die von gleichen Bögen der kleinen Kreise, und von den
gleichen Dreiecksseiten begrenzt werden; diese Zweiecke sind aber congruent,
wie aus der Construction leicht erhellt, also auch die Dreiecke gleich.
\mbox{.\ .\ .\ .\ .\ .\ .}\}

\par\bigskip
